% Chapter 16: The Willmore Conjecture
\let\LocalMacro\relax

\chapter{Willmore Conjecture}

\section{The Problem}

\subsection{Informal}
Imagine inflating a rubber donut and dipping it in soapy water---the film tries to settle into the smoothest shape it can. But not all donut shapes are equally smooth: some are naturally wrinklier, others more rounded. In 1965, Thomas Willmore asked: what is the smoothest a donut-shaped surface can possibly be? He predicted a specific lower limit that no donut shape can ever beat, no matter how you stretch or bend it.

\subsection{Formal}
\begin{conjecture}[Willmore, 1965]
For any immersed torus $T^2 \subset \mathbb{R}^3$,
\[
\mathcal{W}(T^2) \ge 2\pi^2.
\]
Equality holds for the Clifford torus in $S^3$ (stereographically projected to $\mathbb{R}^3$).
\end{conjecture}

\subsection{Historical Context}
Willmore conjectured this in 1965. It became a central problem in geometric analysis, connecting conformal geometry, integrable systems, and algebraic geometry. The conjecture resisted all attempts for 47 years.

\subsection{What Was Known}
Willmore conjectured in 1965 that the minimum energy is $2\pi^2$, achieved by the Clifford torus. Li and Yau proved in 1993 that non-embedded tori have energy at least $2\pi^2$. Ros proved the conjecture for tori of revolution in 2005. But the general case---for all tori, including self-intersecting ones---remained open for 47 years.

\subsection{Why It Matters}
The Willmore energy measures how much a surface deviates from being round---it is the natural measure of ``smoothness'' for curved surfaces. It appears in physics (the elasticity of cell membranes), computer graphics (smoothing 3D models), and biology (the shape of red blood cells). Understanding its minimum value reveals fundamental truths about the geometry of surfaces.

\subsection{Simplified Version}
For spheres (genus 0), the answer is classical: the round sphere minimizes the Willmore energy, achieving the value $4\pi$. The torus (genus 1) is the next-simplest surface, and Willmore conjectured that the Clifford torus---a perfectly symmetric torus living in 4-dimensional space---minimizes the energy among all tori. This was proved by Marques and Neves in 2014.

\section{Mathematical Theory}


\subsection{Prerequisites}
This chapter assumes familiarity with:
\begin{itemize}
\item Differential geometry of surfaces (first/second fundamental forms, mean curvature)
\item Variational calculus and geometric measure theory
\item Basic complex analysis (Riemann surfaces, uniformization)
\item Conformal geometry
\end{itemize}

\subsection{The Willmore Functional and Conformal Invariance}
The Willmore energy is conformally invariant in $\mathbb{R}^3$:
\[
\mathcal{W}(\Sigma) = \int_\Sigma (H^2 - K) \, dA + \int_\Sigma K \, dA = \int_\Sigma (H^2 - K) \, dA + 2\pi\chi(\Sigma)
\]
For a torus ($\chi=0$), $\mathcal{W}(\Sigma) = \int (H^2 - K) \, dA$. The integrand $H^2 - K = \frac{1}{4}(k_1 - k_2)^2 \ge 0$ is conformally invariant.

\subsection{The Clifford Torus}
The Clifford torus in $S^3$ is
\[
T_{\text{Cl}} = \left\{ (z_1, z_2) \in \mathbb{C}^2 : |z_1| = |z_2| = \frac{1}{\sqrt{2}} \right\}.
\]
It is minimal in $S^3$, has constant mean curvature in $\mathbb{R}^3$ (after stereographic projection), and $\mathcal{W}(T_{\text{Cl}}) = 2\pi^2$.

\section{The Solution}

\subsection{What Was Missing Before}

The Willmore conjecture (1965) stated that for any immersed torus $T^2 \subset \mathbb{R}^3$, the Willmore energy $\mathcal{W}(T^2) = \int H^2\,dA \ge 2\pi^2$, with equality only for the Clifford torus (up to conformal transformations). Before the proof, partial results had been obtained. Li and Yau proved in 1993 that non-embedded tori (those with self-intersections) have $\mathcal{W} \ge 2\pi^2$---but this excluded the most interesting case: embedded tori. Ros proved the conjecture for \emph{tori of revolution} in 2005 (surfaces obtained by rotating a circle around an axis). But the general case---for all tori, including embedded ones with complicated topology---remained open. The gap was: nobody could prove the lower bound for \emph{all} immersed tori simultaneously.

\subsection{Why Existing Methods Failed}

The Willmore energy is a \emph{functional} on the space of surfaces---it assigns a number to each torus, and you want to show the minimum is $2\pi^2$. The natural approach is variational: study critical points of the Willmore functional and show the minimum is achieved by the Clifford torus. But this approach fails for a fundamental reason: the space of immersed tori is infinite-dimensional and has complicated topology. You cannot simply ``optimize'' over all tori because:

\begin{enumerate}
\item \textbf{The space is too big.} The space of all immersed tori in $\mathbb{R}^3$ is infinite-dimensional. Direct minimization is impossible---you need to understand the global structure of this space.
\item \textbf{Topology interferes.} Different tori can have different ``topological types'' (different embeddings into $\mathbb{R}^3$), and the Willmore energy is not monotone with respect to any natural ordering. A sequence of tori with decreasing energy might converge to something that is not a torus at all.
\item \textbf{Conformal invariance is a double-edged sword.} The Willmore energy is conformally invariant (invariant under Möbius transformations of $\mathbb{R}^3$), which is a beautiful property but makes the functional ``flat''---it does not distinguish between conformally equivalent tori. This means the energy landscape has large flat directions, making variational methods tricky.
\end{enumerate}

Previous variational approaches (Simon 1993, Bauer and Kuwert 2003) could show that a minimizer exists within certain restricted classes, but they could not handle the full space of immersed tori. The problem was that no one could control the topology of the minimizer: how do you guarantee that a minimizing sequence does not degenerate into something with the wrong genus?

\subsection{What Was Novel}

The breakthrough of Marques and Neves (2014) \cite{marquesneves2014} was the application of \textbf{Almgren--Pitts min-max theory} to the Willmore functional. This theory, originally developed for minimal surfaces (surfaces with $H = 0$), was adapted to the Willmore setting.

The key ideas are:

\begin{enumerate}
\item \textbf{Min-max theory.} Instead of minimizing the Willmore energy directly, you study \emph{sweepouts}---continuous families of tori that fill up a 3-manifold (like $S^3$). For each sweepout, look at the \emph{maximum} Willmore energy among all tori in the family. The min-max value is the infimum of this maximum over all sweepouts. By a compactness argument (using geometric measure theory), this infimum is achieved by a ``min-max torus''---a critical point of the Willmore functional.

\item \textbf{The canonical family.} For any immersed torus $\Sigma$, Marques and Neves construct a \emph{canonical family} of tori in $S^3$ by applying conformal transformations to $\Sigma$. The maximum Willmore energy in this family gives a lower bound on $\mathcal{W}(\Sigma)$. The key is that this family is ``optimal''---no other sweepout can give a better bound.

\item \textbf{The index bound.} Using the min-max theory, Marques and Neves show that any min-max torus with Willmore energy $< 2\pi^2$ must have index $\le 1$ (i.e., it is a local minimum or a saddle point with at most one negative direction). But the only torus in $S^3$ with index $\le 1$ and energy $< 2\pi^2$ is the Clifford torus (which has energy exactly $2\pi^2$). This forces the min-max value to be $\ge 2\pi^2$.

\item \textbf{Free boundary extension.} To handle tori with boundary (which arise in the min-max construction), Marques and Neves use free boundary min-max theory, extending the Almgren--Pitts framework to manifolds with boundary.
\end{enumerate}

The novelty was not just using min-max theory (which had been applied to minimal surfaces since the 1980s) but adapting it to the Willmore functional, which has very different analytic properties. The Willmore functional is fourth-order (involving second derivatives of the embedding), while the area functional is second-order. This required new compactness and regularity theory.

\subsection{Student-Friendly Explanation}

Here is how to think about the proof. The Willmore energy measures how ``round'' a torus is: a perfectly round torus (the Clifford torus) has energy $2\pi^2$, and you want to show no torus can do better.

The natural approach---just minimize the energy---fails because the space of tori is too complicated. Instead, Marques and Neves use a clever indirect strategy:

\begin{enumerate}
\item \textbf{Sweep out $S^3$ with tori.} Imagine filling up the 3-sphere with a continuous family of tori, like inflating a balloon that starts as a tiny Clifford torus and expands to fill all of $S^3$. At some point during the sweepout, the torus must have high Willmore energy---you cannot pass from a tiny torus to a huge one without the energy spiking.

\item \textbf{Find the best sweepout.} Different sweepouts spike at different energies. The min-max value is the best (lowest) spike you can achieve over all sweepouts. Marques and Neves show that this min-max value is exactly $2\pi^2$.

\item \textbf{The Clifford torus is the only possibility.} Using the theory of minimal surfaces in $S^3$, they show that any torus with energy $< 2\pi^2$ must have a very simple structure (low index). But the only such torus is the Clifford torus itself, which has energy exactly $2\pi^2$. So the min-max value cannot be less than $2\pi^2$.

\item \textbf{Conclude.} Since the min-max value is $\ge 2\pi^2$, and the Clifford torus achieves $2\pi^2$, the Willmore conjecture follows: every torus has energy $\ge 2\pi^2$.
\end{enumerate}

The key insight is that min-max theory lets you study the Willmore energy \emph{globally}---you do not need to understand every individual torus, only the structure of sweepouts. The conformal geometry of $S^3$ and the classification of minimal surfaces provide the tools to analyze these sweepouts.

\subsection{Marques and Neves \cite{marquesneves2014}}

\begin{theorem}[Marques-Neves \cite{marquesneves2014}]
The Willmore conjecture is true: $\mathcal{W}(T^2) \ge 2\pi^2$ for any immersed torus $T^2 \subset \mathbb{R}^3$, with equality only for the Clifford torus (up to conformal transformations).
\end{theorem}

Their proof used:
\begin{enumerate}
\item \textbf{Almgren-Pitts min-max theory}: Construct minimal surfaces in $S^3$ via sweepouts.
\item \textbf{The canonical family}: For a given torus $\Sigma$, construct a family of surfaces in $S^3$ by conformal transformations, whose maximum area gives a lower bound on $\mathcal{W}(\Sigma)$.
\item \textbf{The index bound}: Show that any minimal surface in $S^3$ with area $< 2\pi^2$ must be the Clifford torus.
\item \textbf{The free boundary case}: Extend to surfaces with boundary using free boundary min-max.
\end{enumerate}

\begin{highlightbox}
The proof is a tour de force of geometric measure theory: it combines Almgren-Pitts min-max, conformal geometry, and the classification of minimal surfaces in $S^3$ to solve a 47-year-old problem.
\end{highlightbox}

\section{Impact}
\begin{itemize}
\item \textbf{Geometric analysis}: Almgren-Pitts min-max became the standard tool for Willmore-type problems.
\item \textbf{Integrable systems}: Willmore surfaces are related to soliton equations; the Clifford torus corresponds to a finite-gap solution.
\item \textbf{Conformal geometry}: The Willmore energy is the natural energy for conformal immersions.
\item \textbf{Physics}: Willmore energy appears in membrane elasticity (Helfrich energy).
\end{itemize}

\section{Exercises}

\begin{exercise}[Basic]
Compute the Willmore energy of the Clifford torus in $\mathbb{R}^3$ (stereographically projected from $S^3$).
\end{exercise}

\begin{exercise}[Intermediate]
Show that the Willmore energy is conformally invariant in $\mathbb{R}^3$. Why does this fail in higher dimensions?
\end{exercise}

\begin{exercise}[Advanced]
Explain how Almgren-Pitts min-max produces a minimal surface in $S^3$ from a given sweepout of a torus.
\end{exercise}

\section{Timeline}

\begin{tabularx}{\linewidth}{lX}
\toprule
\textbf{Year} & \textbf{Event} \\ \midrule
1965 & Willmore conjectures $2\pi^2$ lower bound for tori \\ 
1993 & Li-Yau proves $\mathcal{W} \ge 2\pi^2$ for non-embedded tori \\ 
2005 & Ros proves Willmore conjecture for tori of revolution \\ 
2012--2014 & Marques--Neves prove Willmore conjecture (preprint 2012, published 2014) \\
2014 & Marques-Neves win Ramanujan Prize \\ \bottomrule
\end{tabularx}

\section{Citations}

\begin{enumerate}
\item Marques, F. C., \& Neves, A. (2014). ``Min-max theory and the Willmore conjecture.'' Annals of Mathematics, 179(2), 683--782.
\item Willmore, T. J. (1965). ``Note on embedded surfaces.'' Anais da Academia Brasileira de Ciências, 37, 475--478.
\item Simon, L. (1993). ``Existence of surfaces minimizing the Willmore functional.'' Communications in Analysis and Geometry, 1(2), 281--326.
\end{enumerate}
